\chapter{Project Objectives}
\label{ch:project_objectives}

The project objectives are formulated according to the SMART method: they are specific, measurable, actionable, relevant, and time-bound. They serve as a roadmap for planning and execution, defining the tangible deliverables and outcomes required within a project’s timeline, budget, and scope.

%The project objectives are split into five main categories: financial, quality, performance, business, and compliance. 

\section{Financial Objectives}
The financial objectives focus on budget management.

\begin{enumerate}
    \item The annual storage cost shall not exceed \EUR{TBD}, per the Den Helder hangar rental agreement.. 
    \item For each landing at Den Helder airport, a fee of \EUR{23.00} shall be paid.
    \item An initial certification fee of \EUR{10,170} applies, followed by a 12-month billing cycle fee of \EUR{11,050} \footnote{\label{ft:EASA_price}\hyperlink{https://eur-lex.europa.eu/eli/reg_impl/2025/2347/oj}{EASA Prices} (accessed on 08.05.2026)}.

\end{enumerate}

\section{Quality Objectives}
The quality objectives define standards for the deliverables.

\begin{enumerate}
    %\item Data shall be downloaded and stored efficiently after each test flight.
    %\item A HUGO configuration shall undergo a thorough inspection of all internal and external systems and the external structure before and after each test flight.
    \item A HUGO configuration shall undergo 50 days of maintenance per year.
    \item The maximum relative difference for verification analysis shall be less than 1\%.
    \item The configuration trade-off shall evaluate the candidate architectures against weighted criteria, with a consistency ratio below 0.1 for the analytic
    hierarchy  (AHP) weighting.
\end{enumerate}

\section{Performance Objectives}
The performance objectives focus on the results and effectiveness of the project.

\begin{enumerate}
    \item The initial HUGO fleet shall comprise of five aircraft to ensure continuous operational availability across 100 operational days per year with 4 flight tests per day.
    \item The selected configuration shall withstand the loads specified in the flight envelope (\autoref{sec:flight_envelope}).
    \item The selected configuration shall satisfy all constraints defined in the matching diagram (\autoref{sec:matching_diagram}).
    \item The selected configuration shall carry sufficient fuel for the nominal mission and 30 minutes of reserve holding time \footnote{\href{https://skybrary.aero/articles/fuel-flight-planning-definitions}{Fuel - Flight Planning Definitions} accessed(20.05.2026)}.
\end{enumerate}

\section{Business Objectives}
The business requirements align with the strategic goals of market entry and client acquisition.
\begin{enumerate}
    \item The Netherlands shall serve as the primary market for initial operations, with Den Helder Airport as the designated base of operations.
    \item HUGO shall be capable of serving customers from at least four market segments: universities and academic laboratories, research centres, aerospace companies, and defence organisations.
    \item The design shall support rapid reconfiguration between client-test campaigns to minimise aircraft-on-ground time between contracts.
\end{enumerate}

\section{Compliance Objectives}
The compliance objectives ensure adherence to regulations and standards.

\begin{enumerate}
    \item Each HUGO configuration shall comply with the environmental performance requirements set out in Annex~III of Regulation (EU) 2018/1139 of the European Parliament and of the Council.
\end{enumerate}